\documentclass{article}
\usepackage{amsthm}
\usepackage{lea-tags}
\newtheorem{claim}{Claim}
\begin{document}
% \leatheorem{label=commented}{This is not a target.}
% \lea{label=unfinished
\lea% Generic definition metadata.
{kind=definition,
 label=syntax_generic,
 context={Use {nested} groups and \{literal\} braces.}
}% The body follows.
{An even natural number is a multiple of two.}

\leatheorem{
 label=syntax_theorem,
 context={Use elementary arithmetic.}% } This brace is commented out.
}% The statement may start after a comment.
{For every natural number $n$, $n+0=n$.}

\lealemma
{label=syntax_lemma}% Statement follows.
{For every natural number $n$, $0+n=n$.}

\leaproposition{label=syntax_proposition}
{For every natural number $n$, $n=n$.}

\leacorollary{label=syntax_corollary}
{For every natural number $n$, $n+1=n+1$.}

\leadefinition{label=syntax_definition}
{The successor of a natural number $n$ is $n+1$.}

\begin{claim}\label{clm:syntax}
\leatheorem
{
 label=syntax_environment,
 context={Use the identity property.}
}
For every natural number $n$, $n=n$.
\end{claim}

\begin{leacode}% Metadata may start on the next line.
{
 label=syntax_code,
 context={Use reflexivity.}
}
theorem syntax_code (n : Nat) : n % 2 = n % 2 := by rfl
-- Literal \lea{label=not_a_tag}, \label{not_a_label}, and { }.
% lea: formalize label=not_a_marker
\end{leacode}
\end{document}
